\documentclass[11pt,a4paper]{article}
\usepackage[margin=1in]{geometry}
\usepackage{amsmath,amssymb,amsthm}
\usepackage{algorithm}
\usepackage{algpseudocode}
\usepackage{booktabs}
\usepackage{hyperref}

\theoremstyle{definition}
\newtheorem{definition}{\textbf{Definition}}[section]
\newtheorem{theorem}{\textbf{Theorem}}[section]
\newtheorem{lemma}[theorem]{\textbf{Lemma}}
\newtheorem{remark}{\textbf{Remark}}[section]

\title{\textbf{Distributed Checkpointing and Exact Resumption Under Elastic Resharding}}
\author{\textbf{Team Scaibu} \\ \small Email: \texttt{jaydeep.9664920749@gmail.com} \\ \small \textbf{Architecture and Core Engine Team}}
\date{\textbf{October 2026}}

\begin{document}
\maketitle

\section{Definitions and Cognitive Mental Models}

\subsection{Formal Mathematical Definition}
\textbf{Definition (Distributed Checkpointing and Elastic Topology Resharding):}
Large-scale pre-training runs lasting months across thousands of GPUs inevitably experience hardware faults, thermal throttling, node preemptions, and silent data corruptions. \textbf{Distributed Checkpointing} (PyTorch DCP; Megatron Distributed Checkpointing) periodically persists complete, consistent distributed training state snapshots $\mathcal{S}_t$ directly from local GPU memory to high-performance parallel file systems (Lustre/GPFS/S3) without requiring consolidation to a single central CPU host.

A complete training checkpoint snapshot $\mathcal{S}_t$ comprises a consistent multi-entity tuple:
{\boldmath
\begin{equation}
\mathcal{S}_t = \left( \theta_{\text{sharded}}^{(t)}, \, \mathcal{M}_{\text{optimizer}}^{(t)}, \, \text{Step}_t, \, \text{LR}_t, \, \text{RNG}_{\text{states}}^{(t)}, \, \text{DataLoaderOffset}_t \right)
\end{equation}
}
Crucially, when resuming training on an elastic cluster whose device count or parallel topology has changed ($N_{\text{save}} \to N_{\text{resume}}$), \textbf{Elastic Resharding} reconstructs the global parameter tensor manifold $\Theta \in \mathbb{R}^P$ from the source shard files and partitions it into $N_{\text{resume}}$ target shards using pre-computed global metadata manifests:
{\boldmath
\begin{equation}
\left\{ \theta_i \right\}_{i=1}^{N_{\text{save}}} \xrightarrow{\text{Global Manifest}} \Theta \xrightarrow{\text{Reshard}} \left\{ \theta_j' \right\}_{j=1}^{N_{\text{resume}}}
\end{equation}
}
Tracking the data loader sample offset $\text{DataLoaderOffset}_t$ guarantees zero-data-leakage and zero-data-skipping: the resumed run continues with exact numerical continuity relative to an uninterrupted run.

\subsection{Expanded Non-Technical Definition (Intuition for Anyone)}
\begin{enumerate}
    \item \textbf{A Quick-Save in an Epic Video Game}: You are on level 90 of a massive game. If your computer crashes or you decide to switch from a laptop to a giant desktop, you don't want to restart at level 1. Distributed checkpointing saves your exact level, character inventory, and location so you continue instantly.
    \item \textbf{Changing the Number of Players Mid-Game (Resharding)}: You saved the game while 8 GPUs were running. The next day, you only have 4 GPUs available, or you upgraded to 16 GPUs. Resharding automatically cuts the saved game files into 4 or 16 equal pieces without corrupting the save file.
    \item \textbf{Exact Data Continuity}: If the checkpoint forgets which book the AI was reading, the AI might re-read the same book 5 times or skip an important encyclopedia entirely! Saving the exact data loader cursor ensures every single sentence is read once and only once.
\end{enumerate}

\subsection{Child-Friendly Mental Model (Physical Metaphor)}
Imagine a team of 8 builders constructing a skyscraper:
\begin{itemize}
    \item \textbf{Unreliable Checkpoint}: Leaving blueprints scattered in the wind. When rain starts, everyone scrambles and forgets which brick goes where.
    \item \textbf{Distributed Checkpointing}: Every 2 hours, all 8 builders snap high-resolution photos of their exact section with a time stamp and save them in a waterproof safe. If tomorrow only 4 builders show up, they look at the master blueprint and continue building smoothly!
\end{itemize}

\section{Core Mathematical Equations and Definitions}

\subsection{Asynchronous Write-Behind Pipeline}
{\boldmath
\begin{equation}
T_{\text{stall}} = \max\left( 0, \, T_{\text{HBM}\to\text{RAM}} - T_{\text{compute}} \right) \ll T_{\text{storage\_IO}}
\end{equation}
}
\begin{itemize}
    \item \textbf{Formal Definition}: Two-stage checkpoint persistence latency under asynchronous write-behind buffering.
    \item \textbf{What It Means}: GPUs copy state to host RAM in $<1$ second and resume training immediately, while background threads flush data to disk.
    \item \textbf{Why It Is Important}: Reduces checkpoint barrier overhead from several minutes down to hundreds of milliseconds.
\end{itemize}

\subsection{Cryptographic Integrity Checksum Verification}
{\boldmath
\begin{equation}
\mathcal{H} = \operatorname{SHA-256}\left( \operatorname{Serialize}(\theta) \parallel \text{Step} \parallel \text{DataLoaderOffset} \right)
\end{equation}
}
\begin{itemize}
    \item \textbf{Formal Definition}: Cryptographic digest validating uncorrupted checkpoint persistence.
    \item \textbf{What It Means}: Protects against silent bit rot, partial writes, and disk corruption.
    \item \textbf{Why It Is Important}: Guarantees cluster does not resume training on corrupted weights.
\end{itemize}

\section{Rigorous Mathematical Analysis Checklist}

\begin{itemize}
    \item \textbf{Notation}: $N_{\text{src}}$ (saving ranks), $N_{\text{dst}}$ (resuming ranks), $P$ (total parameters), $\mathcal{H}$ (checksum).
    \item \textbf{Epistemic Status}: Standardized in PyTorch Distributed Checkpoint (DCP) and Megatron-Core.
    \item \textbf{Dependency Graph}: Local States $\to \text{RAM Snapshot} \to \text{Background Async Disk Flush} \to \text{Manifest Write}$.
    \item \textbf{Dimensions}: Checkpoint byte size $\approx 16P$ bytes for full model state.
    \item \textbf{Regimes}: Production pre-training clusters with mean time between failures (MTBF) $< 48$ hours.
    \item \textbf{Invariants}: Checkpoint resharding perfectly conserves total parameter elements ($\sum |\theta_{\text{dst}}| \equiv \sum |\theta_{\text{src}}|$).
    \item \textbf{Isomorphisms}: Isomorphic to consistent distributed database transaction snapshots and WAL replication.
\end{itemize}

\section{Theoretical Theorems and Formal Proofs}

\begin{theorem}[Conservation of Flat Parameter Manifold Under Partition Resharding]
Let $\Theta \in \mathbb{R}^P$ be a continuous model parameter tensor partitioned across $N_{\text{src}}$ disjoint source shards $\{\theta_i^{\text{src}}\}_{i=1}^{N_{\text{src}}}$. For any target shard count $N_{\text{dst}} \ge 1$ and partition sizes $S_j = \min(P - (j-1)C, C)$ with chunk size $C = \lceil P / N_{\text{dst}} \rceil$:
{\boldmath
\begin{equation}
\bigcup_{j=1}^{N_{\text{dst}}} \theta_j^{\text{dst}} \equiv \bigcup_{i=1}^{N_{\text{src}}} \theta_i^{\text{src}} = \Theta
\end{equation}
}
and total parameter elements satisfy $\sum_{j=1}^{N_{\text{dst}}} |\theta_j^{\text{dst}}| = P$.
\end{theorem}

\begin{proof}
Let index sets $I_i^{\text{src}} = [s_i, e_i)$ define the coordinate slices of source shards such that $\bigcup_{i=1}^{N_{\text{src}}} I_i^{\text{src}} = [0, P)$ and $I_i^{\text{src}} \cap I_k^{\text{src}} = \emptyset$ for $i \neq k$.
The flattened reconstruction operator $\mathcal{R}$ concatenates all source slices:
{\boldmath
\begin{equation}
\Theta[k] = \theta_i^{\text{src}}[k - s_i] \quad \text{for } k \in [s_i, e_i)
\end{equation}
}
Because $\{I_i^{\text{src}}\}$ forms a partition of $[0, P)$, every coordinate $k \in \{0, \dots, P-1\}$ is uniquely defined, giving $\mathcal{R}(\{\theta_i^{\text{src}}\}) = \Theta$.
Similarly, the destination slicing operator $\mathcal{S}$ defines target shards:
{\boldmath
\begin{equation}
\theta_j^{\text{dst}} = \Theta[j \cdot C : \min(P, (j+1) \cdot C)]
\end{equation}
}
The union of destination index intervals $[j \cdot C, \min(P, (j+1)C))$ forms a complete partition of $[0, P)$.
Therefore:
{\boldmath
\begin{equation}
\sum_{j=1}^{N_{\text{dst}}} |\theta_j^{\text{dst}}| = \sum_{j=1}^{N_{\text{dst}}} (\min(P, (j+1)C) - j C) = P
\end{equation}
}
and the reconstructed destination tensor set matches $\Theta$ at every coordinate, proving exact parameter conservation.
\end{proof}

\begin{theorem}[Exact Loss Trajectory Continuity Across Checkpoint Boundary]
Let an uninterrupted training run produce sequence of states $\theta_{t_0}, \theta_{t_0+1}, \dots$ on batch sequence $\mathcal{B}_{t_0}, \mathcal{B}_{t_0+1}, \dots$ with deterministic RNG seed $S$. A training run halted at step $t_0$ and resumed from distributed checkpoint $\mathcal{S}_{t_0}$ reproduces the identical weight trajectory:
{\boldmath
\begin{equation}
\theta_{t_0+1}^{\text{resumed}} \equiv \theta_{t_0+1}^{\text{uninterrupted}}
\end{equation}
}
if and only if $\mathcal{S}_{t_0}$ captures exact DataLoader offset and RNG states.
\end{theorem}

\begin{proof}
By definition of stochastic gradient descent with deterministic optimizer:
{\boldmath
\begin{equation}
\theta_{t+1} = \theta_t - \eta_t \cdot \operatorname{Opt}\left( \nabla \mathcal{L}(\theta_t; \mathcal{B}_t), \, \mathcal{M}_t \right)
\end{equation}
}
The transition map $\theta_t \mapsto \theta_{t+1}$ is a deterministic function of four inputs:
1. Current weights $\theta_t$.
2. Optimizer state $\mathcal{M}_t$ (momentum, variance, step count).
3. Learning rate schedule $\eta_t$.
4. Batch data $\mathcal{B}_t$.
Batch data $\mathcal{B}_t$ is generated by the data loader sampling process driven by the RNG seed and the sample offset:
{\boldmath
\begin{equation}
\mathcal{B}_t = \operatorname{DataLoader}(\operatorname{Offset}_t, \operatorname{RNG}_t)
\end{equation}
}
If $\mathcal{S}_{t_0}$ contains $(\theta_{t_0}, \mathcal{M}_{t_0}, \eta_{t_0}, \operatorname{Offset}_{t_0}, \operatorname{RNG}_{t_0})$, then upon resumption, the inputs to the transition map are identical to those of the uninterrupted run.
By determinism of floating-point arithmetic on identical hardware:
{\boldmath
\begin{equation}
\theta_{t_0+1}^{\text{resumed}} = \theta_{t_0} - \eta_{t_0} \operatorname{Opt}\left( \nabla \mathcal{L}(\theta_{t_0}; \mathcal{B}_{t_0}), \mathcal{M}_{t_0} \right) = \theta_{t_0+1}^{\text{uninterrupted}}
\end{equation}
}
By induction, the entire subsequent training trajectory is identical, completing the proof.
\end{proof}

\section{Foundational Architectural Questions and Epistemic Meta-Questions}

\subsection{Architectural Question 1: Checkpoint Frequency Optimization}
\textbf{Question:} How often should checkpoints be saved during large-scale training?\\
\textbf{Answer:} Using Young's formula for optimal checkpoint interval under Poisson failure rates:
{\boldmath
\begin{equation}
\Delta t_{\text{optimal}} = \sqrt{2 \cdot T_{\text{save}} \cdot \text{MTBF}}
\end{equation}
}
where $T_{\text{save}}$ is the checkpoint write duration and $\text{MTBF}$ is the Mean Time Between Failures of the cluster. For a 1,000-GPU cluster with $\text{MTBF} \approx 24$ hours and $T_{\text{save}} \approx 10$ minutes, optimal checkpoint frequency is every $2.8$ hours.

\subsection{Architectural Question 2: Silent Data Corruption (SDC) Protection}
\textbf{Question:} Why are checksums mandatory for distributed checkpoints?\\
\textbf{Answer:} Cosmic rays, PCIe bus errors, and failing storage controllers occasionally cause Silent Data Corruption (single-bit flips). A single bit flip in an FP16 weight matrix can turn a valid parameter into `NaN` or `1e38`, causing instantaneous model collapse. Cryptographic SHA-256 validation catches corrupted shards before weights are loaded into VRAM.

\section{Algorithmic Implementation Specification}

\begin{algorithm}[H]
\caption{Distributed Checkpoint Serialization and Resharding}
\begin{algorithmic}[1]
\Require Source shards $\{\mathbf{\theta}_i^{\text{src}}\}_{i=1}^{N_{\text{src}}}$, step $t$, data offset $\mathcal{O}$, target shard count $N_{\text{dst}}$
\Ensure Target shards $\{\mathbf{\theta}_j^{\text{dst}}\}_{j=1}^{N_{\text{dst}}}$, metadata manifest, checksum $\mathcal{H}$, total elements $P$
\State Flatten source shards: $\mathbf{\Theta} \gets \operatorname{Concat}(\mathbf{\theta}_1^{\text{src}}, \dots, \mathbf{\theta}_{N_{\text{src}}}^{\text{src}})$
\State $P \gets |\mathbf{\Theta}|$
\State $C \gets \lceil P / N_{\text{dst}} \rceil$
\State Allocate $\{\mathbf{\theta}_j^{\text{dst}}\} \gets [\,]$
\For{$j = 0$ \textbf{to} $N_{\text{dst}} - 1$}
    \State $\text{start} \gets j \cdot C, \quad \text{end} \gets \min(P, (j + 1) \cdot C)$
    \State Append $\mathbf{\Theta}[\text{start} : \text{end}]$ to $\{\mathbf{\theta}_j^{\text{dst}}\}$
\EndFor
\State Compute digest $\mathcal{H} \gets \operatorname{SHA-256}(\mathbf{\Theta}[0 : \min(1000, P)] \parallel t \parallel \mathcal{O})$
\State Construct manifest $\mathcal{M} \gets \{ \text{"step"}: t, \, \text{"offset"}: \mathcal{O}, \, \text{"elements"}: P, \, \text{"shards"}: N_{\text{dst}} \}$
\State \Return $(\{\mathbf{\theta}_j^{\text{dst}}\}, \mathcal{M}, \mathcal{H}, P)$
\end{algorithmic}
\end{algorithm}

\section{Big-$\Theta$ Complexity Analysis}

\begin{itemize}
    \item \textbf{Sharding Flattening Pass}: $\Theta(P)$ linear memory movement.
    \item \textbf{Target Slicing Operation}: $\Theta(P)$ partition assignment.
    \item \textbf{Integrity Hashing Time}: $\Theta(\min(P, 1000))$ constant bounded digest.
    \item \textbf{Total Time Complexity}: $\Theta(P)$ linear operations.
    \item \textbf{Space Complexity}: $\Theta(P)$ temporary staging buffer.
\end{itemize}

\section{Single Canonical Repository Reference}

The complete deterministic scalar implementation, verification suite, and unit test benchmarks are published at:
\begin{center}
\url{https://github.com/Chief-Strategist-J/llm-obs-infra}
\end{center}

\end{document}
